% Apunts de teoria en PDF · 4t ESO · Unitat 5: Sistemes d'equacions i inequacions
% Mateix estil que els apunts de Batxillerat (franja i caixes de teoria en blau, exemples en roig), però més visuals,
% curts i amb dibuixos, per a l'alumnat de 3r ESO. Es compila amb pdflatex dues vegades (per als enllaços de l'índex),
% amb logo-fentmates.pdf (_eines/logo/pdf/horitzontal-clar.pdf) al costat.
\documentclass[11pt,a4paper]{article}
\usepackage[utf8]{inputenc}
\usepackage[T1]{fontenc}
\usepackage[catalan]{babel}
\usepackage[sfdefault,book,lining]{FiraSans}
\usepackage[italic,noplusnominus,noequal]{mathastext}
\linespread{1.18}
\usepackage[a4paper,top=2.4cm,bottom=2.2cm,left=1.6cm,right=1.6cm,headheight=1.1cm,headsep=0.5cm,footskip=1cm]{geometry}
\usepackage{amsmath,amssymb,cancel}
\usepackage{icomma}
\usepackage{graphicx}
\usepackage[table]{xcolor}
\usepackage{tikz}\usetikzlibrary{calc,positioning,arrows.meta,decorations.pathreplacing,shapes.geometric}
\usepackage{fontawesome5}
\usepackage{ccicons}
\usepackage{fancyhdr}
\usepackage{enumitem}
\usepackage{array}
\usepackage{needspace}
\usepackage[most]{tcolorbox}
\usepackage{titlesec}
\usepackage[hidelinks]{hyperref}
\hypersetup{pdftitle={Apunts de teoria. Unitat 5: Sistemes d'equacions i inequacions. 4t ESO},pdfsubject={Matemàtiques A 4t ESO},pdfkeywords={sistemes lineals, mètode gràfic, substitució, igualació, reducció, parèntesis i denominadors, sistemes no lineals, problemes, inequacions de primer grau, inequacions de segon grau, sistemes d'inequacions}}

% colors de la web (bloc d'Àlgebra: blau)
\definecolor{fmnavy}{HTML}{1B2748}
\definecolor{fmblue}{HTML}{3D6AA0}
\definecolor{fmbluel}{HTML}{4E81BD}
\definecolor{fmyellow}{HTML}{FFC94D}
\definecolor{fmgold}{HTML}{E9A919}
\definecolor{fmaccent}{HTML}{FF7A50}
\definecolor{fmdef}{HTML}{1D4F73}\definecolor{fmdefbg}{HTML}{EEF4F9}
\definecolor{fmex}{HTML}{A3203A}\definecolor{fmexbg}{HTML}{FBF0F2}
\definecolor{fmsf}{HTML}{C2410C}
\definecolor{fmok}{HTML}{15803D}
\definecolor{fmmuted}{HTML}{5D6685}
\definecolor{fmsurface}{HTML}{EEF1F8}
\definecolor{fmA}{HTML}{2F6FD6}\definecolor{fmB}{HTML}{EA6A10}\definecolor{fmR}{HTML}{E11D48}\definecolor{fmK}{HTML}{16A34A}\definecolor{fmV}{HTML}{7C3AED}

% lletres
\newcommand{\comic}{\fontfamily{ComicNeue-TLF}\selectfont}
\newcommand{\heavy}{\fontseries{eb}\selectfont}
\newcommand{\semi}{\fontseries{sb}\selectfont}
\DeclareUnicodeCharacter{2212}{-}

\setlength{\parindent}{0pt}
\setlength{\parskip}{7pt}
\setlength{\lineskip}{4pt}\setlength{\lineskiplimit}{4pt}
\setlength{\emergencystretch}{3em}
\setlist{itemsep=3pt,topsep=3pt,leftmargin=1.4em}
\setcounter{secnumdepth}{1}
\renewcommand{\thesection}{5.\arabic{section}}
\titleformat{\section}[hang]{\color{fmnavy}\fontsize{23}{27}\heavy\selectfont}{\tcbox[on line,colback=fmblue,colframe=fmblue,arc=3mm,boxsep=2pt,left=6pt,right=6pt,top=3pt,bottom=3pt]{\color{white}\thesection}}{0.45em}{}[\vspace{5pt}{\color{fmblue}\titlerule[2pt]}]
\titlespacing*{\section}{0pt}{4pt}{6pt}
\titleformat{\subsection}{\color{fmblue}\Large\semi}{}{0pt}{\faCaretRight\ }
\titlespacing*{\subsection}{0pt}{8pt}{2pt}

% frase de cada apartat, «a mà»
\newcommand{\frase}[1]{\par\vspace{-2pt}{\comic\Large\color{fmbluel}\faCommentDots[regular]\ \ #1\par}\vspace{2pt}}
\newcommand{\seccio}[3]{\Needspace{16\baselineskip}\section{#2}\hypertarget{#1}{}\frase{#3}}
% notes a mà dins dels exemples
\newcommand{\nota}[1]{{\comic\small\color{fmex!85!black}#1}}
\newcommand{\notab}[1]{{\comic\small\color{fmblue}#1}}
% bé / malament
\newcommand{\be}[1]{{\color{fmok}\faCheckCircle}\ \ensuremath{#1}}
\newcommand{\mal}[1]{{\color{fmR}\faTimesCircle}\ {\color{fmR}\ensuremath{#1}}}
% pas numerat
\newcommand{\pas}[1]{\tikz[baseline=(p.base)]\node[circle,fill=fmblue,text=white,inner sep=1.2pt,minimum size=1.25em,font=\small\heavy](p){#1};}
% etiqueta de color
\newcommand{\chip}[2]{\tcbox[on line,colback=#1!14,colframe=#1!60,arc=2mm,boxsep=1pt,left=5pt,right=5pt,top=1pt,bottom=1pt]{#2}}
\newcommand{\fmkey}[1]{\colorbox{fmyellow!55}{\strut #1}}
\newcommand{\cA}[1]{{\color{fmA}#1}}\newcommand{\cB}[1]{{\color{fmB}#1}}\newcommand{\cK}[1]{{\color{fmK}#1}}\newcommand{\cR}[1]{{\color{fmR}#1}}\newcommand{\cV}[1]{{\color{fmV}#1}}

% caixes
\tcbset{fmbase/.style={enhanced,arc=3mm,boxrule=1.2pt,left=4mm,right=4mm,top=2mm,bottom=2.5mm,before skip=8pt,after skip=8pt,fonttitle=\heavy\large,before upper=\raggedright}}
\newtcolorbox{idea}[1][Idea clau]{fmbase,colback=fmdefbg,colframe=fmdef,coltitle=white,title={\faKey\ \ #1}}
\newtcolorbox{exemple}[1][Exemple]{fmbase,colback=fmexbg,colframe=fmex,colbacktitle=fmexbg,coltitle=fmex,title={\faPencil*\ \ #1},titlerule=0.5pt,titlerule style={fmex!40,dashed}}
\newtcolorbox{compte}[1][Compte!]{fmbase,colback=fmsf!6,colframe=fmsf,coltitle=white,fonttitle=\comic\bfseries\Large,title={\faExclamationTriangle\ \ #1}}
\newtcolorbox{truc}[1][Truc]{fmbase,colback=fmyellow!20,colframe=fmgold,colbacktitle=fmyellow,coltitle=fmnavy,fonttitle=\comic\bfseries\Large,title={\faLightbulb[regular]\ \ #1}}

% dibuix de la targeta de la unitat, en blanc
\newcommand{\fmcard}{\begin{tikzpicture}[x=0.062cm,y=-0.062cm,white,line width=3pt,line cap=round]
  \begin{scope}[transparency group,opacity=0.15]
  \node[anchor=base west,inner sep=0pt,font=\fontsize{26}{26}\selectfont] at (14,44) {$\left\{\begin{array}{l}x+y=7\\ x-y=1\end{array}\right.$};
  \draw[line width=3pt] (150,110) -- (235,30) (150,40) -- (235,100);
  \end{scope}
\end{tikzpicture}}
\newcommand{\fmbrand}{\raisebox{-0.25\height}{\includegraphics[height=0.5cm]{logo-fentmates.pdf}}}

% regla de Ruffini: #1 = a, #2 = coeficients, #3 = fila del mig (sense el primer), #4 = fila de baix, #5 = nombre de coeficients
\newcommand{\ruffini}[5]{\begin{tikzpicture}[x=1.25cm,y=0.85cm,every node/.style={font=\large}]
  \foreach \v [count=\i] in {#2} \node (t\i) at (\i,0) {$\v$};
  \foreach \v [count=\i] in {#3} {\pgfmathtruncatemacro\j{\i+1}\node (m\j) at (\j,-1) {$\v$};}
  \foreach \v [count=\i] in {#4} \node (b\i) at (\i,-2.1) {$\v$};
  \node[text=fmB] at (0,-1) {$#1$};
  \draw[line width=1.1pt,fmnavy] (0.45,0.5)--(0.45,-1.5) (0.45,-1.5)--(#5+0.5,-1.5);
  \node[draw=fmR,line width=1.3pt,rounded corners=2mm,minimum size=0.75cm] at (#5,-2.1) {};
  \draw[-{Latex[length=2.2mm]},fmK,line width=1.1pt] (t1.south) -- (b1.north);
  \pgfmathtruncatemacro\last{#5-1}
  \foreach \i in {1,...,\last} {\pgfmathtruncatemacro\j{\i+1}
    \draw[-{Latex[length=2.2mm]},fmB,line width=1.1pt,shorten >=1pt] (b\i.north east) -- (m\j.south west);
    \draw[-{Latex[length=2.2mm]},fmA!70,line width=1pt,dashed,shorten >=2pt] ([xshift=8pt]t\j.south) -- ([xshift=8pt]b\j.north);}
\end{tikzpicture}}

% capçalera i peu
\pagestyle{fancy}\fancyhf{}
\fancyhead[L]{\small\fmbrand}
\fancyhead[R]{\small\color{fmmuted}4t ESO · Unitat 5: Sistemes d'equacions i inequacions}
\fancyfoot[L]{\scriptsize\color{fmmuted}\href{https://creativecommons.org/licenses/by-nc-sa/4.0/deed.ca}{\ccbyncsa}\ \ FentMates · fentmates.org · CC BY-NC-SA 4.0}
\fancyfoot[R]{\small\color{fmmuted}\thepage}
\renewcommand{\headrulewidth}{0.4pt}\renewcommand{\headrule}{\color{fmmuted!40}\hrule height \headrulewidth}
\fancypagestyle{primera}{\fancyhf{}\renewcommand{\headrulewidth}{0pt}
  \fancyfoot[L]{\scriptsize\color{fmmuted}\href{https://creativecommons.org/licenses/by-nc-sa/4.0/deed.ca}{\ccbyncsa}\ \ FentMates · fentmates.org · CC BY-NC-SA 4.0}
  \fancyfoot[R]{\small\color{fmmuted}\thepage}}

% targeta de l'índex
\newcommand{\itile}[4]{\begin{tcolorbox}[enhanced,colback=white,colframe=fmblue!35,arc=3mm,boxrule=1pt,boxsep=0pt,left=6pt,right=6pt,top=5pt,bottom=5pt,before skip=0pt,after skip=7pt]
  \hyperlink{#1}{\begin{minipage}[c]{0.9cm}\centering{\color{fmblue}\LARGE #3}\end{minipage}\hspace{6pt}%
  \begin{minipage}[c]{0.78\linewidth}{\color{fmblue}\heavy #2}\\[-2pt]{\color{fmnavy}\semi\large #4}\end{minipage}}\end{tcolorbox}}

% taula amb vores arrodonides: el contingut es retalla amb la forma (sense @{} a les vores ni entre columnes)
\newsavebox\tabclipbox
\newcommand{\taulaclip}[3][3mm]{\sbox\tabclipbox{#3}\begin{tikzpicture}[baseline=(tcn.base)]
  \begin{scope}\clip[rounded corners=#1] (0,-\dp\tabclipbox) rectangle (\wd\tabclipbox,\ht\tabclipbox);
  \node[inner sep=0pt,outer sep=0pt,anchor=base west] (tcn) at (0,0) {\usebox\tabclipbox};\end{scope}
  \draw[#2,rounded corners=#1] (0,-\dp\tabclipbox) rectangle (\wd\tabclipbox,\ht\tabclipbox);
\end{tikzpicture}}
\newcommand{\taularod}[3][fmblue]{\taulaclip{#1,line width=1.2pt}{\renewcommand{\arraystretch}{1.5}\begin{tabular}{#2}#3\end{tabular}}}
% balança
\newcommand{\balanca}[3]{\begin{tikzpicture}[scale=#3,transform shape,every node/.style={anchor=center}]
  \fill[fmnavy] (-0.45,-1.9) rectangle (0.45,-1.75);
  \fill[fmnavy!80] (-0.12,-1.75) rectangle (0.12,0);
  \fill[fmgold] (0,0.05) circle (0.13);
  \draw[fmnavy,line width=2.4pt,line cap=round] (-2,0) -- (2,0);
  \foreach \s in {-1,1} {
    \draw[fmmuted,line width=0.8pt] (\s*2,0) -- (\s*2-0.75,-1.05) (\s*2,0) -- (\s*2+0.75,-1.05);
    \fill[fmyellow!70,draw=fmgold,line width=1pt] (\s*2-0.95,-1.05) arc[start angle=180,end angle=360,x radius=0.95,y radius=0.32] -- cycle;}
  \node[font=\Large,fill=fmA!15,draw=fmA,rounded corners=1.5mm,inner sep=3pt] at (-2,-0.68) {$#1$};
  \node[font=\Large,fill=fmB!15,draw=fmB,rounded corners=1.5mm,inner sep=3pt] at (2,-0.68) {$#2$};
\end{tikzpicture}}
% eixos amb quadrícula (dins d'un tikzpicture): xmin, xmax, ymin, ymax
\newcommand{\eixos}[4]{\draw[fmmuted!22,line width=0.4pt] (#1,#3) grid (#2,#4);
  \draw[-{Latex[length=2mm]},fmmuted,line width=0.8pt] (#1-0.2,0)--(#2+0.4,0) node[right,font=\small]{$x$};
  \draw[-{Latex[length=2mm]},fmmuted,line width=0.8pt] (0,#3-0.2)--(0,#4+0.4) node[above,font=\small]{$y$};}
% taula de traduccions (diccionari) amb les vores arrodonides
\newcommand{\dic}[2]{#1 & {\large$#2$}\\}
\newcommand{\dicratio}{0.62}
\newcommand{\dichead}{\begin{tabular}{@{}>{\raggedright}p{\dimexpr\dicratio\linewidth-12pt}>{\centering\arraybackslash}p{\dimexpr\linewidth-\dicratio\linewidth-12pt}@{}}\faComment[regular]\ Ho diem així\dots & \dots\ així\end{tabular}}
\newenvironment{dicbox}{\begin{tcolorbox}[enhanced,arc=3mm,boxrule=1.2pt,colframe=fmblue,colback=white,coltitle=white,fonttitle=\heavy,
  title={\dichead},left=0pt,right=0pt,top=0pt,bottom=1.5mm,boxsep=0pt,lefttitle=6pt,righttitle=6pt,before skip=0pt,after skip=0pt]\rowcolors{1}{white}{fmsurface}%
  \begin{tabular}{@{\hspace{6pt}}>{\raggedright}p{\dimexpr\dicratio\linewidth-12pt}>{\centering\arraybackslash}p{\dimexpr\linewidth-\dicratio\linewidth-12pt}@{\hspace{6pt}}}}
  {\end{tabular}\end{tcolorbox}}
% resum
\newcommand{\rtile}[3]{\begin{tcolorbox}[enhanced,colback=white,colframe=fmblue!45,arc=3mm,boxrule=1pt,left=3mm,right=3mm,top=2.5mm,bottom=2.5mm,fontupper=\large,
  title={\tcbox[on line,colback=fmblue,colframe=fmblue,arc=2mm,boxsep=0pt,left=4pt,right=4pt,top=1pt,bottom=1pt]{\color{white}\small\heavy #1}\ \ #2},coltitle=fmnavy,colbacktitle=fmdefbg,fonttitle=\semi,toptitle=0.5mm,bottomtitle=0.5mm]#3\end{tcolorbox}}
% portada: #1 subtítol a mà
\newcommand{\portada}[1]{\thispagestyle{primera}
\begin{tikzpicture}[remember picture,overlay]
  \fill[fmblue] (current page.north west) rectangle ([yshift=-9.4cm]current page.north east);
  \begin{scope}
    \clip (current page.north west) rectangle ([yshift=-9.4cm]current page.north east);
    \node[anchor=north west,inner sep=0pt] at ([xshift=6.6cm,yshift=-1.3cm]current page.north west) {\fmcard};
  \end{scope}
  \node[anchor=north west,inner sep=0pt] at ([xshift=1.6cm,yshift=-1.1cm]current page.north west) {\tcbox[colback=white,colframe=white,arc=2.5mm,boxsep=0pt,left=7pt,right=7pt,top=5pt,bottom=5pt]{\includegraphics[height=0.75cm]{logo-fentmates.pdf}}};
\end{tikzpicture}
\vspace*{0.55cm}
\tcbox[colback=white!15!fmblue,colframe=white!15!fmblue,arc=2.5mm,boxsep=0pt,left=6pt,right=6pt,top=3pt,bottom=3pt]{\color{white}\small\bfseries Apunts de teoria · 4t ESO · Matemàtiques A}

\vspace{2pt}
{\raggedright\color{white}\fontsize{30}{36}\heavy\selectfont{\color{fmyellow}Unitat 5}\\[2pt]Sistemes d'equacions i inequacions\par}
\vspace{8pt}
{\color{white}\comic\Large #1\par}
\vspace*{\dimexpr 7.4cm-\pagetotal\relax}
}
\newcommand{\llegenda}{\par\vspace{16pt}
{\color{fmnavy}\heavy\Large\faPalette\ \ Com es llegeixen aquests apunts}\par\smallskip
\begin{tcbraster}[raster columns=4,raster equal height,raster column skip=8pt,raster before skip=12pt]
\tcbox[colback=fmdefbg,colframe=fmdef,arc=3mm,boxrule=1.2pt,halign=center,width=\linewidth]{{\color{fmdef}\Large\faKey}\\[2pt]{\color{fmdef}\heavy Idea clau}\\[-1pt]{\small el que cal saber}}
\tcbox[colback=fmexbg,colframe=fmex,arc=3mm,boxrule=1.2pt,halign=center,width=\linewidth]{{\color{fmex}\Large\faPencil*}\\[2pt]{\color{fmex}\heavy Exemple}\\[-1pt]{\small pas a pas}}
\tcbox[colback=fmsf!6,colframe=fmsf,arc=3mm,boxrule=1.2pt,halign=center,width=\linewidth]{{\color{fmsf}\Large\faExclamationTriangle}\\[2pt]{\color{fmsf}\heavy Compte!}\\[-1pt]{\small errors típics}}
\tcbox[colback=fmyellow!20,colframe=fmgold,arc=3mm,boxrule=1.2pt,halign=center,width=\linewidth]{{\color{fmgold!80!black}\Large\faLightbulb[regular]}\\[2pt]{\color{fmnavy}\heavy Truc}\\[-1pt]{\small per a recordar-ho}}
\end{tcbraster}}
\newcommand{\resumtitol}{\clearpage
\hypertarget{resum}{}
{\color{fmnavy}\fontsize{23}{27}\heavy\selectfont\faClipboardList\ \ En resum}\par\vspace{5pt}{\color{fmblue}\titlerule[2pt]}
\frase{Tota la unitat en una pàgina.}}
% els 5 passos per a resoldre problemes: #2 = text del pas 2
\newcommand{\passosproblema}[2]{\begin{center}
\begin{tikzpicture}[s/.style={draw=fmblue,fill=fmdefbg,rounded corners=3mm,line width=1.2pt,minimum width=2.75cm,minimum height=2.45cm,align=center,font=\small},
  ic/.style={font=\LARGE,text=fmblue}]
  \foreach \i/\ic/\t in {1/\faBookOpen/{\textbf{Llegim}\\i entenem},2/\faQuestionCircle/{#1},3/\faEquals/{#2},4/\faCogs/{\textbf{Resolem}},5/\faCheckDouble/{\textbf{Comprovem}\\i responem}} {
    \node[s] (s\i) at (3.3*\i,0) {};
    \node[ic] at ([yshift=-0.5cm]s\i.north) {\ic};
    \node[align=center,font=\small,anchor=north] at ([yshift=-0.95cm]s\i.north) {\t};
    \node[circle,fill=fmblue,text=white,font=\small\heavy,inner sep=1.5pt,minimum size=1.3em] at (s\i.north west) {\i};}
  \foreach \i [evaluate=\i as \j using int(\i+1)] in {1,...,4} \draw[-{Latex[length=3mm]},fmmuted,line width=1.5pt] (s\i) -- (s\j);
\end{tikzpicture}
\end{center}}

\newcommand{\lo}{\mathopen{]}}\newcommand{\ro}{\mathclose{[}}
% punts ple i buit, nombres dels eixos i taula de valors
\newcommand{\pple}[2][fmblue]{\fill[#1] #2 circle (3.6pt);}
\newcommand{\nums}[4]{\foreach \k in {#1,...,#2} {\ifnum\k=0\else\node[below,font=\scriptsize,text=fmmuted] at (\k,0) {$\k$};\fi}
  \foreach \k in {#3,...,#4} {\ifnum\k=0\else\node[left,font=\scriptsize,text=fmmuted] at (0,\k) {$\k$};\fi}}
\newcommand{\tv}[3][fmblue]{\taulaclip{#1,line width=1.2pt}{\renewcommand{\arraystretch}{1.45}%
  \begin{tabular}{>{\columncolor{#1}\color{white}\heavy}c*{12}{w{c}{0.8cm}}}#2\\\rowcolor{#1!8}\cellcolor{#1}#3\end{tabular}}}
% etiqueta d'una recta o corba
\newcommand{\etq}[3]{\node[fill=white,inner sep=1.5pt,text=#1,font=\small] at #2 {#3};}


\def\intvdot#1#2#3{\ifx#2o\filldraw[fill=white,draw=#3,line width=1.3pt] (#1,0) circle (3.3pt);\fi\ifx#2c\fill[#3] (#1,0) circle (3.5pt);\fi}
\newcommand{\intv}[5]{\begin{tikzpicture}[x=0.4cm,baseline=-3pt]
  \draw[fmmuted,line width=0.8pt] (-6,0)--(7,0);
  \foreach \k in {-5,...,6} \draw[fmmuted,line width=0.6pt] (\k,-0.08)--(\k,0.08);
  \ifx#3a\def\xa{-6}\else\def\xa{#1}\fi
  \ifx#4a\def\xb{7}\else\def\xb{#2}\fi
  \ifx#4a\draw[#5,line width=2.6pt,-{Latex[length=2.6mm]}] (\xa,0)--(\xb,0);\else\ifx#3a\draw[#5,line width=2.6pt,{Latex[length=2.6mm]}-] (\xa,0)--(\xb,0);\else\draw[#5,line width=2.6pt] (\xa,0)--(\xb,0);\fi\fi
  \intvdot{#1}{#3}{#5}\intvdot{#2}{#4}{#5}
  \ifx#3a\else\node[below=2pt,font=\scriptsize] at (#1,0) {$#1$};\fi
  \ifx#4a\else\node[below=2pt,font=\scriptsize] at (#2,0) {$#2$};\fi
\end{tikzpicture}}
\newcommand{\sis}[2]{\left\{\begin{array}{@{}l@{}}#1\\[2pt]#2\end{array}\right.}
\newcommand{\passos}[1]{\begin{tcolorbox}[fmbase,colback=fmdefbg,colframe=fmdef,coltitle=white,title={\faListOl\ \ Com es resolen}]#1\end{tcolorbox}}


\begin{document}
\portada{Dues equacions alhora, i desigualtats que es resolen amb intervals.}

{\color{fmnavy}\heavy\Large\faMapSigns\ \ Què aprendràs?}\par\smallskip
\begin{minipage}[t]{0.485\linewidth}
\itile{sec1}{5.1}{\faChartLine}{Sistemes lineals. Mètode gràfic}
\itile{sec2}{5.2}{\faIcon{exchange-alt}}{Substitució, igualació, reducció}
\itile{sec3}{5.3}{\faIcon{layer-group}}{Parèntesis i denominadors}
\itile{sec4}{5.4}{\faIcon{bezier-curve}}{Sistemes no lineals}
\itile{sec5}{5.5}{\faTasks}{Problemes amb sistemes}
\end{minipage}\hfill
\begin{minipage}[t]{0.485\linewidth}
\itile{sec6}{5.6}{\faAngleRight}{Inequacions de primer grau}
\itile{sec7}{5.7}{\faChartArea}{Inequacions de segon grau}
\itile{sec8}{5.8}{\faRulerHorizontal}{Sistemes d'inequacions}
\itile{resum}{\phantom{5.8}}{\faClipboardList}{En resum}
\end{minipage}
\llegenda

% ============================================================ 5.1
\clearpage
\seccio{sec1}{Sistemes lineals. Mètode gràfic}{Cada equació és una recta; la solució, on es tallen.}

\vfill
\begin{minipage}[c]{0.5\linewidth}
\begin{idea}[Sistema de dues equacions lineals]
{\centering\Large$\sis{a_1x+b_1y=c_1}{a_2x+b_2y=c_2}$\par}\smallskip
Una \textbf{solució} és un parell $(x,y)$ que compleix les dues alhora: el \textbf{punt de tall} de les dues rectes.
\end{idea}
\end{minipage}\hfill
\begin{minipage}[c]{0.47\linewidth}\centering
\begin{tikzpicture}[scale=0.48]
  \eixos{-1}{5}{-1}{7}
  \nums{1}{4}{1}{7}
  \draw[fmA,line width=1.6pt] (0,7) -- (3.5,0); \draw[fmB,line width=1.6pt] (1,-1) -- (5,3);
  \fill[fmR] (3,1) circle (6pt); \node[right=3pt,font=\small\heavy] at (3.1,1.3) {$(3,\,1)$};
  \node[fill=white,inner sep=1pt,font=\small,text=fmA] at (2.4,4.4) {$2x+y=7$}; \node[fill=white,inner sep=1pt,font=\small,text=fmB] at (4.2,3.4) {$x-y=2$};
\end{tikzpicture}
\end{minipage}

\vfill
\begin{center}\small
\taularod{>{\heavy}l<{\hspace{4pt}}>{\hspace{4pt}}c<{\hspace{4pt}}>{\hspace{4pt}}l<{\hspace{4pt}}>{\hspace{4pt}}c<{\hspace{4pt}}}{
\rowcolor{fmblue}\color{white}Rectes & \color{white}\heavy Solucions & \color{white}\heavy Sistema & \color{white}\heavy Coeficients\\
Es tallen & una & compatible determinat (SCD) & $\frac{a_1}{a_2}\neq\frac{b_1}{b_2}$\rule[-8pt]{0pt}{22pt}\\
\rowcolor{fmsurface}Paral·leles & cap & incompatible (SI) & $\frac{a_1}{a_2}=\frac{b_1}{b_2}\neq\frac{c_1}{c_2}$\rule[-8pt]{0pt}{22pt}\\
Coincidents & infinites & compatible indeterminat (SCI) & $\frac{a_1}{a_2}=\frac{b_1}{b_2}=\frac{c_1}{c_2}$\rule[-8pt]{0pt}{22pt}
}
\end{center}

\vfill
\begin{minipage}[t]{0.5\linewidth}
\passos{\pas{1} Una taula amb dos punts per a cada recta.\\[3pt]
\pas{2} Les dues rectes en els mateixos eixos.\\[3pt]
\pas{3} Llegim el \textbf{punt de tall}.}
\end{minipage}\hfill
\begin{minipage}[t]{0.47\linewidth}
\begin{compte}[Poc precís]
Si la solució no és entera, el dibuix només l'aproxima. Normalment fem servir els mètodes algebraics (5.2).
\end{compte}
\end{minipage}

\vfill
\begin{exemple}[Exemples]
$\sis{2x+y=7}{x-y=2}$: \ $y=7-2x$ per $(0,7)$ i $(3,1)$; \ $y=x-2$ per $(2,0)$ i $(3,1)$ \ $\to$ \ \boxed{(3,\,1)}\\[5pt]
$\sis{x-y=1}{2x-2y=5}$: \ $\dfrac12=\dfrac{-1}{-2}\neq\dfrac15$ \ $\to$ paral·leles, \textbf{incompatible} \qquad
$\sis{x+3y=2}{2x+6y=4}$: \ la mateixa recta, \textbf{SCI}
\end{exemple}

\vfill

\vfill
% ============================================================ 5.2
\clearpage
\seccio{sec2}{Substitució, igualació i reducció}{Eliminem una incògnita i queda una equació de primer grau.}

\vfill
\begin{tcbraster}[raster columns=3,raster equal height,raster column skip=8pt]
\begin{tcolorbox}[fmbase,colback=fmA!8,colframe=fmA,coltitle=white,title={Substitució}]
Aïllem una incògnita i la substituïm en l'altra. Còmode si té coeficient $\pm1$.
\end{tcolorbox}
\begin{tcolorbox}[fmbase,colback=fmB!8,colframe=fmB,coltitle=white,title={Igualació}]
Aïllem la \textbf{mateixa} incògnita en les dues i igualem.
\end{tcolorbox}
\begin{tcolorbox}[fmbase,colback=fmK!8,colframe=fmK,coltitle=white,title={Reducció}]
Multipliquem perquè una incògnita tinga coeficients oposats i sumem.
\end{tcolorbox}
\end{tcbraster}

\vfill
\begin{exemple}[Exemple: \ el mateix sistema $\sis{x+2y=8}{3x-y=3}$ pels tres mètodes]
\begin{minipage}[t]{0.32\linewidth}
\textbf{\cA{Substitució}}\\[2pt]
$x=8-2y$\\
$3(8-2y)-y=3$\\
$-7y=-21\to y=3$\\
$x=8-6=2$
\end{minipage}\hfill
\begin{minipage}[t]{0.32\linewidth}
\textbf{\cB{Igualació}}\\[2pt]
$y=\dfrac{8-x}{2}$, \ $y=3x-3$\\[2pt]
$8-x=6x-6$\\
$x=2$, \ $y=3$
\end{minipage}\hfill
\begin{minipage}[t]{0.32\linewidth}
\textbf{\cK{Reducció}}\\[2pt]
2a per 2: \ $6x-2y=6$\\
sumem: \ $7x=14$\\
$x=2$, \ $y=3$
\end{minipage}\par\medskip
{\centering Solució: \ \boxed{x=2,\ y=3} \ \nota{comprova-la en les dues}\par}
\end{exemple}

\vfill
\begin{tcbraster}[raster columns=2,raster equal height,raster column skip=10pt]
\begin{tcolorbox}[fmbase,colback=fmK!8,colframe=fmK,coltitle=white,title={Arribem a $0=0$}]
Compatible \textbf{indeterminat}: infinites solucions.
\end{tcolorbox}
\begin{tcolorbox}[fmbase,colback=fmR!6,colframe=fmR,coltitle=white,title={Arribem a $0=5$}]
Igualtat impossible: \textbf{incompatible}.
\end{tcolorbox}
\end{tcbraster}

\vfill
\begin{exemple}[Exemple: \ amb parèntesis i denominadors]
$\sis{\frac{x+1}{2}+\frac y3=2}{2(x-y)=x-3}$ \ $\to$ \ primer la forma $ax+by=c$: \ $\sis{3x+2y=9}{x-2y=-3}$ \ $\to$ \ sumant, $4x=6$: \ \boxed{x=\tfrac32,\ y=\tfrac94}
\end{exemple}

\vfill
\begin{exemple}[Exemples]
$\sis{y=3x-5}{2x+y=10}$: \ $5x=15$, \ $(3,\,4)$ \qquad
$\sis{x=2y+1}{x=5-2y}$: \ $y=1$, \ $(3,\,1)$ \qquad
$\sis{3x+2y=12}{5x-2y=4}$: \ $8x=16$, \ $(2,\,3)$
\end{exemple}

\vfill

\vfill
% ============================================================ 5.3
\clearpage
\seccio{sec3}{Sistemes amb parèntesis i denominadors}{Primer els deixem nets; després, com sempre.}

\vfill
\begin{idea}
{\centering\Large Transformem cada equació en la forma \ \fmkey{$ax+by=c$}\par}\smallskip
Igual que en les equacions de primer grau: llevem \textbf{denominadors} (pel m.c.m.), llevem \textbf{parèntesis}, transposem i reduïm. Si tots els coeficients tenen un divisor comú, \textbf{simplifiquem}.
\end{idea}

\vfill
\begin{exemple}[Exemple: \ $\sis{2(x+1)-y=3(y-1)+5}{\dfrac{x+2}{3}+\dfrac{y-1}{2}=\dfrac52}$]
\begin{minipage}[t]{0.52\linewidth}
\chip{fmA}{1a} \ $2x+2-y=3y-3+5$\\[3pt]
\hspace*{2.2em}$2x-4y=0\ \to\ \cA{x-2y=0}$ \ \nota{entre 2}\\[8pt]
\chip{fmB}{2a} \ \nota{m.c.m.$=6$:} \ $2(x+2)+3(y-1)=15$\\[3pt]
\hspace*{2.2em}$2x+4+3y-3=15\ \to\ \cB{2x+3y=14}$
\end{minipage}\hfill
\begin{minipage}[t]{0.45\linewidth}
Ja és un sistema net:\\[3pt]
{\large$\sis{\cA{x-2y=0}}{\cB{2x+3y=14}}$}\\[6pt]
Substitució: \ $x=2y\ \to\ 4y+3y=14$\\[3pt]
$\boxed{y=2}$, \quad $\boxed{x=4}$
\end{minipage}
\end{exemple}

\vfill
\subsection{Quin mètode trie?}
\providecommand{\metode}{}\renewcommand{\metode}[4]{\begin{tcolorbox}[fmbase,colback=fmdefbg,colframe=fmdef,coltitle=white,title={#1\ \ #2}]
#3\par\smallskip{\centering#4\par}
\end{tcolorbox}}
\begin{tcbraster}[raster columns=2,raster equal height=rows,raster column skip=10pt]
\metode{\faIcon{exchange-alt}}{Substitució}{Una incògnita té coeficient $1$ o $-1$, o ja està aïllada.}{$\sis{y=2x-3}{3x+2y=8}$}
\metode{\faEquals}{Igualació}{La mateixa incògnita està aïllada (o és fàcil d'aïllar) en les dues.}{$\sis{x=2y+1}{x=7-y}$}
\metode{\faCompress}{Reducció}{Coeficients iguals, oposats o múltiples; o cap coeficient 1.}{$\sis{3x+2y=8}{5x-2y=8}$}
\metode{\faChartLine}{Gràfic}{Per a veure el sistema i de quin tipus és. Només és exacte si la solució és entera.}{$\sis{x+y=5}{2x-y=1}$}
\end{tcbraster}
{\comic\color{fmmuted}\centering Qualsevol mètode dona la mateixa solució. Si vols, comprova-la amb un altre.\par}

\vfill
% ============================================================ 5.4
\clearpage
\seccio{sec4}{Sistemes no lineals}{Quadrats, productes\dots: corbes que es tallen.}

\vfill
\begin{minipage}[c]{0.55\linewidth}
\begin{idea}[Sistemes no lineals]
Alguna equació no és de primer grau ($x^2$, $y^2$, $xy$\dots).\\[3pt]
\textbf{Substitució}: aïllem de la lineal i substituïm. Sol quedar una de segon grau.\\[3pt]
\textbf{Reducció}: si $x^2$ i $y^2$ apareixen igual en les dues.\\[3pt]
Pot haver-hi \textbf{diverses solucions}: per cada valor d'una incògnita, calculem l'altra.
\end{idea}
\end{minipage}\hfill
\begin{minipage}[c]{0.42\linewidth}\centering
\begin{tikzpicture}[scale=0.55]
  \eixos{-3}{3}{-2}{4}
  \draw[fmA,line width=1.6pt,domain=-2.3:2.3,samples=40] plot (\x,{\x*\x-1});
  \draw[fmB,line width=1.6pt] (-2.8,-1.8) -- (2.8,3.8);
  \fill[fmR] (2,3) circle (5pt) node[right,font=\small] {$(2,3)$}; \fill[fmR] (-1,0) circle (5pt) node[above left,font=\small] {$(-1,0)$};
\end{tikzpicture}\\
{\comic\small\color{fmmuted}$y=x^2-1$ \ i \ $y=x+1$}
\end{minipage}

\vfill
\begin{exemple}[Exemple: \ $\sis{x-y=1}{x^2+y^2=13}$]
$x=1+y\ \to\ (1+y)^2+y^2=13\ \to\ y^2+y-6=0$ \qquad $y=2\to x=3$; \ $y=-3\to x=-2$ \qquad \boxed{(3,2)\ \text{i}\ (-2,-3)}
\end{exemple}

\vfill
\begin{exemple}[Exemple: \ $\sis{x+y=7}{xy=12}$]
$y=7-x\ \to\ x(7-x)=12\ \to\ x^2-7x+12=0$: \ $x=3,\ x=4$ \qquad \boxed{(3,4)\ \text{i}\ (4,3)}
\end{exemple}

\vfill
\begin{exemple}[Exemple: \ $\sis{x^2+y^2=25}{x^2-y^2=7}$ \ \ per reducció]
Sumem: \ $2x^2=32\ \to\ x=\pm4$; \ $y^2=9\ \to\ y=\pm3$ \qquad quatre solucions: \ $(\pm4,\,\pm3)$
\end{exemple}

\vfill
\begin{exemple}[Exemple: \ $\sis{y=x^2-1}{y=x+1}$ \ \ paràbola i recta]
Igualem: \ $x^2-1=x+1\ \to\ x^2-x-2=0$: \ $x=2,\ x=-1$ \qquad \boxed{(2,3)\ \text{i}\ (-1,0)} \ \nota{els punts del dibuix}
\end{exemple}

\vfill
\begin{compte}[Calcula sempre l'altra incògnita]
Cada valor de $x$ dona el seu $y$: no barreges les parelles. \ $\sis{x+y=0}{xy=1}$ dona $-x^2=1$: \textbf{no té solució}.
\end{compte}

\vfill

\vfill
% ============================================================ 5.5
\clearpage
\seccio{sec5}{Problemes amb sistemes}{Cada condició de l'enunciat, una equació.}

\passosproblema{\textbf{Triem}\\$x$ i $y$}{\textbf{Plantegem}\\el sistema}

\vfill
\begin{exemple}[Exemple: \ 30 vehicles entre cotxes i motos, i 100 rodes]
\begin{minipage}[c]{0.5\linewidth}
$x$ = cotxes, \ $y$ = motos: \ $\sis{x+y=30}{4x+2y=100}$
\end{minipage}\hfill
\begin{minipage}[c]{0.47\linewidth}
Restem el doble de la 1a: \ $2x=40$\\
\boxed{20\text{ cotxes i }10\text{ motos}}
\end{minipage}
\end{exemple}

\vfill
\begin{exemple}[Exemple: \ un rectangle de 32 cm de perímetre i 60 cm$^2$ d'àrea]
\begin{minipage}[c]{0.5\linewidth}
$\sis{2x+2y=32}{xy=60}$ \ $\to$ \ $y=16-x$
\end{minipage}\hfill
\begin{minipage}[c]{0.47\linewidth}
$x^2-16x+60=0$: \ $x=10,\ x=6$\\
\boxed{10\text{ cm i }6\text{ cm}}
\end{minipage}
\end{exemple}

\vfill
\begin{exemple}[Exemple: \ dos nombres es diferencien en 3 i la suma dels quadrats és 117]
$\sis{x-y=3}{x^2+y^2=117}$ \ $\to$ \ $y^2+3y-54=0$: \ $y=6\to x=9$; \ $y=-9\to x=-6$ \qquad \boxed{9\text{ i }6,\ \text{o}\ -6\text{ i }-9}
\end{exemple}

\vfill
\begin{exemple}[Exemples]
250 animals entre conills i gallines, 800 potes: \ $\sis{x+y=250}{4x+2y=800}$ \ $\to$ \ 150 conills i 100 gallines\\[5pt]
Xifres que sumen 9; en intercanviar-les, el nombre augmenta 27: \ $\sis{x+y=9}{(10y+x)-(10x+y)=27}$ \ $\to$ \ \textbf{36}\\[5pt]
Rectangle de diagonal 10 cm i àrea 48 cm$^2$: \ $\sis{x^2+y^2=100}{xy=48}$ \ $\to$ \ 8 cm i 6 cm
\end{exemple}

\vfill
\begin{compte}[Interpreta les solucions]
Descarta les que no tenen sentit (longituds negatives, persones no enteres\dots) i respon amb una frase.
\end{compte}

\vfill

\vfill
% ============================================================ 5.6
\clearpage
\seccio{sec6}{Inequacions de primer grau}{Com una equació, però amb una regla especial.}

\vfill
\begin{idea}[Inequació]
Una desigualtat ($<$, $\le$, $>$, $\ge$) amb una incògnita. Les solucions solen ser \textbf{infinites}: les escrivim amb \textbf{intervals}.
\end{idea}

\vfill
\begin{tcbraster}[raster columns=2,raster equal height,raster column skip=10pt]
\begin{tcolorbox}[fmbase,colback=fmK!8,colframe=fmK,coltitle=white,title={\faCheck\ \ Es manté}]
Sumar o restar, i multiplicar o dividir per un nombre \textbf{positiu}.
\end{tcolorbox}
\begin{tcolorbox}[fmbase,colback=fmR!6,colframe=fmR,coltitle=white,title={\faIcon{exchange-alt}\ \ Canvia de sentit}]
Multiplicar o dividir per un \textbf{negatiu}:\par\smallskip
{\centering$-2x<6\ \to\ x\cR{>}-3$\par}
\end{tcolorbox}
\end{tcbraster}

\vfill
\begin{tcbraster}[raster columns=2,raster equal height,raster column skip=10pt]
\begin{exemple}[Exemple: \ $3x-2\le x+4$]
$2x\le6\ \to\ x\le3$ \qquad $\boxed{\lo-\infty,\,3]}$
\end{exemple}
\begin{exemple}[Exemple: \ $-2x+1<7$]
$-2x<6$. \ Entre $-2$, \textbf{girem}: \ $x>-3$ \qquad $\boxed{\lo-3,+\infty\ro}$
\end{exemple}
\end{tcbraster}

\vfill
\begin{exemple}[Exemple: \ $\dfrac{x-1}{2}-\dfrac{x+2}{3}\ge1$]
Per 6 (positiu): \ $3(x-1)-2(x+2)\ge6\ \to\ 3x-3-2x-4\ge6\ \to\ x\ge13$ \qquad $\boxed{[13,+\infty\ro}$
\end{exemple}

\vfill
\begin{exemple}[Exemples]
$5x+3>2x-6\to x>-3$ \qquad $4-3x\ge10\to x\le-2$ \qquad $2(x-1)<5x+4\to x>-2$ \qquad $\frac x2-1\le\frac x3\to x\le6$\\[4pt]
Amb parèntesis: \ $3(x-2)-2(1-x)>4x+1\ \to\ 5x-8>4x+1\ \to\ x>9$: \ $\lo9,+\infty\ro$
\end{exemple}

\vfill
\begin{exemple}[Exemple: \ sistema \ $2x-1>3$ \ i \ $x+4\le10$]
\begin{minipage}[c]{0.5\linewidth}
Resolem cada una: \ $x>2$ \ i \ $x\le6$\\[3pt]
Les solucions són les \textbf{comunes} (intersecció):\\[3pt]
$\boxed{\lo2,\,6]}$
\end{minipage}\hfill
\begin{minipage}[c]{0.47\linewidth}\centering
\begin{tikzpicture}[x=0.5cm]
  \draw[fmmuted,line width=0.8pt] (-1,0) -- (9,0); \foreach \k in {0,...,8} \draw[fmmuted] (\k,-0.07) -- (\k,0.07);
  \draw[fmA,line width=2.4pt,-{Latex[length=2.4mm]}] (2,0.35) -- (9,0.35); \filldraw[fill=white,draw=fmA,line width=1.2pt] (2,0.35) circle (3pt);
  \draw[fmB,line width=2.4pt,{Latex[length=2.4mm]}-] (-1,-0.35) -- (6,-0.35); \fill[fmB] (6,-0.35) circle (3.2pt);
  \fill[fmK!40,opacity=0.6] (2,-0.15) rectangle (6,0.15);
  \node[below=6pt,font=\scriptsize] at (2,-0.35) {2}; \node[below=6pt,font=\scriptsize] at (6,-0.35) {6};
\end{tikzpicture}
\end{minipage}
\end{exemple}

\vfill
\begin{minipage}[t]{0.48\linewidth}
\begin{exemple}[Exemple: \ $-3<2x+1\le7$]
Restem 1 als tres membres: \ $-4<2x\le6$\\[3pt]
Dividim entre 2: \ $-2<x\le3$ \ $\to$ \ $\boxed{\lo-2,\,3]}$
\end{exemple}
\end{minipage}\hfill
\begin{minipage}[t]{0.49\linewidth}
\begin{exemple}[Problema: \ tarifa A: 20 € + 0,05 €/min; \ tarifa B: 0,10 €/min. Quan és millor A?]
$20+0{,}05x<0{,}10x\ \to\ 20<0{,}05x$\\[3pt]
$x>400$: \ \boxed{\text{més de 400 min}}
\end{exemple}
\end{minipage}

\vfill

\vfill
% ============================================================ 5.7
\clearpage
\seccio{sec7}{Inequacions de segon grau o superior}{Arrels, intervals i taula de signes.}

\vfill
\begin{minipage}[c]{0.52\linewidth}
\passos{\pas{1} Tot a un membre: \ $P(x)>0$ \ o \ $P(x)<0$ \ (o amb $\ge$, $\le$).\\[3pt]
\pas{2} Resolem l'equació $P(x)=0$.\\[3pt]
\pas{3} Situem les solucions en la recta real i vegem el signe en cada tros substituint en $P(x)$.\\[3pt]
\pas{4} Triem els intervals. Les arrels, només amb $\le$ o $\ge$.}
\end{minipage}\hfill
\begin{minipage}[c]{0.45\linewidth}\centering
\begin{tikzpicture}[x=0.7cm,y=0.55cm]
  \draw[-{Latex[length=2mm]},fmmuted,line width=0.8pt] (-0.5,0) -- (5.5,0) node[right,font=\scriptsize] {$x$};
  \draw[fmblue,line width=1.6pt,domain=0.9:4.1,samples=40] plot (\x,{(\x-2)*(\x-3)*1.5});
  \draw[fmK,line width=4pt,opacity=0.6] (-0.5,0) -- (2,0) (3,0) -- (5.5,0);
  \draw[fmR,line width=4pt,opacity=0.6] (2,0) -- (3,0);
  \node[below,font=\scriptsize] at (2,-0.1) {2}; \node[below,font=\scriptsize] at (3,-0.1) {3};
  \node[font=\small\heavy,text=fmK!80!black] at (0.8,1.2) {$+$}; \node[font=\small\heavy,text=fmR] at (2.5,-1) {$-$}; \node[font=\small\heavy,text=fmK!80!black] at (4.2,1.2) {$+$};
\end{tikzpicture}\\
{\comic\small\color{fmmuted}$a>0$: negativa \textbf{entre} les arrels}
\end{minipage}


\vfill
\begin{exemple}[Exemple: \ $x^3-x^2-6x\ge0$]
\begin{minipage}[c]{0.62\linewidth}\centering\small
\renewcommand{\arraystretch}{1.25}
\taularod[fmex]{>{\heavy}c*{7}{>{\hspace{2pt}}c<{\hspace{2pt}}}}{
\rowcolor{fmex} & \color{white}$\lo-\infty,-2\ro$ & \color{white}$-2$ & \color{white}$\lo-2,0\ro$ & \color{white}$0$ & \color{white}$\lo0,3\ro$ & \color{white}$3$ & \color{white}$\lo3,+\infty\ro$\\
$x+2$ & $-$ & 0 & $+$ & $+$ & $+$ & $+$ & $+$\\
\rowcolor{fmexbg}$x$ & $-$ & $-$ & $-$ & 0 & $+$ & $+$ & $+$\\
$x-3$ & $-$ & $-$ & $-$ & $-$ & $-$ & 0 & $+$\\
\rowcolor{fmexbg}$P(x)$ & \cR{$-$} & 0 & \cK{$+$} & 0 & \cR{$-$} & 0 & \cK{$+$}
}
\end{minipage}\hfill
\begin{minipage}[c]{0.35\linewidth}
$x(x+2)(x-3)\ge0$\\[4pt]
$\boxed{[-2,\,0]\cup[3,+\infty\ro}$
\end{minipage}
\end{exemple}

\vfill
\begin{exemple}[Exemples]
$-x^2+x+2\ge0\ \to\ x^2-x-2\le0$ \ \nota{per $-1$, girem}: \ $[-1,\,2]$ \qquad $x^2\ge9$: \ $\lo-\infty,-3]\cup[3,+\infty\ro$\\[4pt]
$(x+1)^2>0$: \ $\mathbb{R}-\{-1\}$ \qquad $x^2+1<0$: \ \textbf{cap solució} \qquad $x^3-4x>0$: \ $\lo-2,0\ro\cup\lo2,+\infty\ro$
\end{exemple}

\vfill
\begin{exemple}[Exemple: \ $x^2-x-6\le0$, \ amb la recta]
\begin{minipage}[c]{0.45\linewidth}
\pas{1} $x^2-x-6=0\ \to\ x=-2,\ x=3$\\[3pt]
\pas{2} Un valor de cada tros:\\ $P(-3)=6$, \ $P(0)=-6$, \ $P(4)=6$\\[3pt]
\pas{3} Volem $\le0$: \ $\boxed{[-2,\,3]}$ \ \nota{amb les arrels}
\end{minipage}\hfill
\begin{minipage}[c]{0.52\linewidth}\centering
\begin{tikzpicture}[x=0.75cm]
  \draw[-{Latex[length=2mm]},fmmuted,line width=0.8pt] (-5,0) -- (5.5,0);
  \draw[fmR,line width=4pt,opacity=0.6] (-2,0) -- (3,0);
  \fill[fmnavy] (-2,0) circle (3pt) node[below=3pt,font=\scriptsize] {$-2$}; \fill[fmnavy] (3,0) circle (3pt) node[below=3pt,font=\scriptsize] {$3$};
  \node[font=\heavy,text=fmK!80!black] at (-3.6,0.45) {$+$}; \node[font=\heavy,text=fmR] at (0.5,0.45) {$-$}; \node[font=\heavy,text=fmK!80!black] at (4.3,0.45) {$+$};
  \node[font=\scriptsize,text=fmmuted] at (-3.6,-0.45) {$x=-3$}; \node[font=\scriptsize,text=fmmuted] at (0.5,-0.45) {$x=0$}; \node[font=\scriptsize,text=fmmuted] at (4.3,-0.45) {$x=4$};
\end{tikzpicture}
\end{minipage}
\end{exemple}

\vfill
\begin{tcbraster}[raster columns=2,raster equal height,raster column skip=10pt]
\begin{exemple}[Exemple: \ $(x+1)(x-3)<2x-3$]
$x^2-2x-3<2x-3\ \to\ x^2-4x<0$\\[3pt]
$x(x-4)<0$: \ arrels 0 i 4; \ $P(1)=-3<0$\\[3pt]
$\boxed{\lo0,\,4\ro}$
\end{exemple}
\begin{exemple}[Problema: \ rectangle de 20 m de perímetre. Quina base fa que l'àrea passe de 21 m$^2$?]
$x(10-x)>21\ \to\ x^2-10x+21<0$\\[3pt]
Arrels 3 i 7: \ \boxed{\text{entre 3 i 7 m}}
\end{exemple}
\end{tcbraster}

\vfill

\vfill
% ============================================================ 5.8
\clearpage
\seccio{sec8}{Sistemes d'inequacions}{Resolem cada una i ens quedem amb el que tenen en comú.}

\vfill
\begin{idea}[Com es resolen]
Resolem \textbf{cada inequació per separat} (5.6 i 5.7). La solució del sistema és la \textbf{intersecció}. Si no tenen cap punt en comú, el sistema \textbf{no té solució}.
\end{idea}

\vfill
\begin{exemple}[Exemple: \ $\sis{2x-1<5}{x^2-4\ge0}$]
\begin{minipage}[c]{0.55\linewidth}
\begin{tabular}{@{}r@{\ \ }l@{}}
$x<3$ & \intv{-6}{3}{a}{o}{fmA}\\[4pt]
$x^2\ge4$ & \begin{tikzpicture}[x=0.4cm,baseline=-3pt]
  \draw[fmmuted,line width=0.8pt] (-6,0)--(7,0); \foreach \k in {-5,...,6} \draw[fmmuted,line width=0.6pt] (\k,-0.08)--(\k,0.08);
  \draw[fmB,line width=2.6pt,{Latex[length=2.6mm]}-] (-6,0) -- (-2,0); \draw[fmB,line width=2.6pt,-{Latex[length=2.6mm]}] (2,0) -- (7,0);
  \fill[fmB] (-2,0) circle (3.5pt) node[below=2pt,font=\scriptsize,text=black] {$-2$}; \fill[fmB] (2,0) circle (3.5pt) node[below=2pt,font=\scriptsize,text=black] {$2$};
\end{tikzpicture}\\[4pt]
\cK{sistema} & \begin{tikzpicture}[x=0.4cm,baseline=-3pt]
  \draw[fmmuted,line width=0.8pt] (-6,0)--(7,0); \foreach \k in {-5,...,6} \draw[fmmuted,line width=0.6pt] (\k,-0.08)--(\k,0.08);
  \draw[fmK,line width=2.6pt,{Latex[length=2.6mm]}-] (-6,0) -- (-2,0); \draw[fmK,line width=2.6pt] (2,0) -- (3,0);
  \fill[fmK] (-2,0) circle (3.5pt); \fill[fmK] (2,0) circle (3.5pt); \filldraw[fill=white,draw=fmK,line width=1.3pt] (3,0) circle (3.3pt);
\end{tikzpicture}
\end{tabular}
\end{minipage}\hfill
\begin{minipage}[c]{0.42\linewidth}
$2x<6\to x<3$\\[3pt]
$x^2-4\ge0\to\lo-\infty,-2]\cup[2,+\infty\ro$\\[5pt]
\boxed{\lo-\infty,-2]\cup[2,\,3\ro}\\[3pt]
\nota{el 3 no: no compleix la primera}
\end{minipage}
\end{exemple}

\vfill
\begin{exemple}[Exemple: \ $\sis{x-2>0}{x+1<0}$]
$x>2$ \ i \ $x<-1$: \ els dos intervals no tenen cap punt en comú \ $\to$ \ \textbf{no té solució}
\end{exemple}

\vfill
\begin{exemple}[Exemples]
$\sis{2x-3\le0}{3x+1>0}$: \ $\left]-\frac13,\,\frac32\right]$ \qquad
$\sis{x^2-x-6<0}{x\ge0}$: \ $\lo-2,3\ro\cap[0,+\infty\ro=[0,\,3\ro$ \qquad
$\sis{x+1\ge0}{x<-3}$: \ \textbf{cap}
\end{exemple}

\vfill
\begin{truc}[Una recta damunt de l'altra]
Dibuixa la solució de cada inequació en una recta, amb la mateixa escala: el que coincideix en totes és la solució.
\end{truc}

\vfill

\vfill
% ============================================================ resum
\resumtitol
\begin{tcbraster}[raster columns=2,raster equal height=rows,raster column skip=14pt,raster row skip=34pt,raster before skip=18pt]
\rtile{5.1}{Mètode gràfic}{Punt de tall de les rectes.\\[1ex] Es tallen: SCD. Paral·leles: SI.\\[1ex] Coincidents: SCI.}
\rtile{5.2}{Mètodes algebraics}{Substitució, igualació o reducció.\\[1ex] $0=0$: SCI. \ $0=5$: SI.}
\rtile{5.3}{Parèntesis i denominadors}{Primer, cada equació a la forma\par\vspace{1ex}{\centering$ax+by=c$\par}}
\rtile{5.4}{No lineals}{Substitució (o reducció).\\[1ex] Calcula sempre l'altra incògnita.}
\rtile{5.5}{Problemes}{Una condició, una equació.\\[1ex] Interpreta les solucions.}
\rtile{5.6}{Inequacions 1r grau}{Per un negatiu, la desigualtat gira.\\[1ex] Solució: un interval.}
\rtile{5.7}{Grau 2 o més}{Arrels i signe de cada tros.}
\rtile{5.8}{Sistemes}{Cada una per separat; la solució és la intersecció.}
\end{tcbraster}
\end{document}
